\scanpage{182}
% archive sheet number: 181

\[\underset{\unsure{m|\mathfrak m}}{A}\qquad\deleted{\unsure{M}}\]

$X,Y$ deux modules sur $A$
\[H^0_m(X)\xrightarrow{\sim}DH^0_m(Y)\]

$(X,Y)\unsure{\longleftarrow|\longrightarrow}(X',Y')$ consiste en
\[\begin{aligned}X&\longrightarrow X'\\Y&\longleftarrow Y'\end{aligned}\]

tel que le diagramme
\diagram{\begin{tikzcd}[ampersand replacement=\&]
H^0(X) \arrow[r] \arrow[d,"\sim"'] \& H^0(X') \arrow[d,"\sim"] \\
DH^0(Y) \arrow[r,dash] \& DH^0(Y')
\end{tikzcd}}
soit \uuline{commutatif}

Si on veut inclure les groupes $G_m$ (\deleted{\unsure{dans les}} parmi les algébriques (avec aussi $\mu_p,\mathbb Z/\ell\mathbb Z$) et $\mathbb Z$ parmi les groupes formels (avec aussi $\mathbb Z/p\mathbb Z,\mathbb Z/\ell\mathbb Z$), alors la \unsure{prés.} est la $=$, en prenant comme nouvel anneau topologique
\[\mathbb Z\times A\]
et le $H^0$ étant relatif aux idéaux $m\unsure{=|:}(F,V)$ de $A$ et les idéaux non nuls quelconques de $\mathbb Z$, la dualité
